% =============================================================================
% Pre-Projeto de Mestrado — PPGCC / DCC / UFMG
% Área: Criptografia Aplicada / Sistemas Distribuídos / Blockchain
% Orientador: Prof. Jeroen van de Graaf
% =============================================================================
\documentclass[12pt,a4paper]{article}

% --- Pacotes ---
\usepackage[utf8]{inputenc}
\usepackage[T1]{fontenc}
\usepackage[brazil]{babel}
\usepackage{amsmath,amssymb,amsfonts}
\usepackage{hyperref}
\usepackage{graphicx}
\usepackage{geometry}
\usepackage{setspace}
\usepackage{cite}
\usepackage{lipsum}
\usepackage{enumitem}
\usepackage{fancyhdr}

\geometry{top=3cm, bottom=2.5cm, left=3cm, right=2.5cm}
\onehalfspacing

% --- Metadados do PDF ---
\hypersetup{
  pdftitle={Pre-Projeto de Mestrado — Criptografia Aplicada a Sistemas Distribuídos e Blockchain},
  pdfauthor={Arthur de Freitas Abeilice},
  pdfsubject={Mestrado em Ciência da Computação},
  pdfkeywords={zero-knowledge proofs, blockchain, cross-chain, storage proofs, consensus proofs, criptografia aplicada}
}

% --- Cabeçalho ---
\pagestyle{fancy}
\fancyhf{}
\fancyhead[L]{\small Pre-Projeto de Mestrado — PPGCC/UFMG}
\fancyhead[R]{\small Abeilice, A.F.}
\fancyfoot[C]{\thepage}
\renewcommand{\headrulewidth}{0.4pt}

% =============================================================================
\begin{document}

% --- Capa ---
\begin{titlepage}
  \centering
  \vspace*{3cm}

  {\LARGE \textbf{Pré-Projeto de Mestrado}\par}
  \vspace{1.5cm}

  {\Large \textbf{Provas de Conhecimento-Zero para Interoperabilidade}\\
           \textbf{Descentralizada entre Blockchains}\par}
  \vspace{0.5cm}
  \textit{(Zero-Knowledge Proofs for Decentralized Blockchain Interoperability)}

  \vspace{2cm}

  {\large \textbf{Arthur de Freitas Abeilice}\par}
  \vspace{0.5cm}
  {\large Orientador: Prof. Dr. Jeroen van de Graaf\par}

  \vspace{2cm}

  {\large Programa de Pós-Graduação em Ciência da Computação\\
           Departamento de Ciência da Computação\\
           Universidade Federal de Minas Gerais\par}

  \vspace{1cm}
  {\large Belo Horizonte — MG\\
           Julho de 2026\par}

  \vfill
\end{titlepage}

% --- Resumo ---
\begin{abstract}
\lipsum[1-2]

\noindent\textbf{Palavras-chave:} Provas de conhecimento-zero, interoperabilidade entre blockchains, provas de armazenamento, provas de consenso, sistemas distribuídos.
\end{abstract}

\newpage

% --- Abstract (English) ---
\renewcommand{\abstractname}{Abstract}
\begin{abstract}
\lipsum[1-2]

\noindent\textbf{Keywords:} Zero-knowledge proofs, blockchain interoperability, storage proofs, consensus proofs, distributed systems.
\end{abstract}

\newpage

\tableofcontents
\newpage

% =============================================================================
% 1. INTRODUÇÃO
% =============================================================================
\section{Introdução}

\lipsum[3-5]

\subsection{Contextualização}

\lipsum[6-7]

\subsection{Motivação}

\lipsum[8-9]

\subsection{Objetivos}

\subsubsection{Objetivo Geral}
\lipsum[10]

\subsubsection{Objetivos Específicos}
\begin{enumerate}[label=(\roman*)]
  \item \lipsum[11][1-2]
  \item \lipsum[11][3-4]
  \item \lipsum[11][5-6]
  \item \lipsum[12][1-2]
\end{enumerate}

\subsection{Estrutura do Documento}
\lipsum[13]

% =============================================================================
% 2. FUNDAMENTAÇÃO TEÓRICA
% =============================================================================
\section{Fundamentação Teórica e Revisão da Literatura}

\lipsum[14]

\subsection{Blockchain e Sistemas Distribuídos}

A arquitetura fundamental de blockchains como sistemas distribuídos é estabelecida pelo
\textit{Ethereum Yellow Paper} \cite{ethereum_yellow_paper}, que define o modelo de
estado baseado em Merkle Patricia Trie. O problema de consenso em sistemas
distribuídos é extensivamente tratado por Bano et al. \cite{consensus_age_of_blockchains}.

\lipsum[15-16]

\subsection{Provas de Conhecimento-Zero}

O protocolo PLONK \cite{aztec_plonk} estabelece uma base para argumentos de
conhecimento não-interativos. Extensões como PLOOKUP \cite{aztec_plookup}
introduzem tabelas de consulta (\textit{lookup tables}) para otimização de
circuitos. O protocolo Aztec \cite{aztec_protocol_whitepaper} aplica ZK para
privacidade em rollups.

\paragraph{Sistemas baseados em STARKs.} A plataforma Miden introduz avanços em
STARKs com foco em máquinas virtuais \cite{miden_arithmetization_oriented_encryption},
adição de conhecimento-zero a STARKs \cite{miden_adding_zk_to_starks}, e
assinaturas baseadas em funções de hash RPO \cite{miden_stark_signatures_rpo}.

\lipsum[17]

\subsection{Interoperabilidade entre Blockchains}

A ponte zkBridge \cite{zkbridge} demonstra a viabilidade de pontes
\textit{trustless} utilizando provas ZK. O protocolo Zendoo
\cite{zendoo} propõe sidechains desacopladas com verificabilidade via
zk-SNARKs. A fundação formal de rollups é estabelecida em \cite{formal_foundation_blockchain_rollups}.

\lipsum[18]

\subsection{Provas de Armazenamento e Consenso}

Herodotus introduz provas de armazenamento históricas e \textit{multichain}
\cite{historical_multichain_storage_proofs}, com o comitê DALOC para
disponibilidade de dados \cite{herodotus_daloc}. O protocolo Bankai
\cite{bankai} propõe provas de consenso para clientes leves. Uma
sistematização do conhecimento sobre clientes leves baseados em SNARKs
é apresentada em \cite{snark_light_clients_sok}.

\lipsum[19]

\subsection{Privacidade e Anonimato em Blockchain}

O Semaphore \cite{semaphore_whitepaper} fornece sinalização anônima
com ZK. UniRep \cite{unirep_pse_report} estende o conceito para
reputação universal. PLUME \cite{plume_ecdsa_nullifier} introduz
\textit{nullifiers} baseados em ECDSA. WAKU RLN \cite{waku_rln_relay}
aplica \textit{rate-limiting nullifiers} para comunicação descentralizada.

\lipsum[20]

\subsection{Ferramentas e Técnicas de Verificação}

A Veridise \cite{veridise_underconstrained_circuits} analisa circuitos
sub-restritos em provas ZK. O Halo \cite{halo_recursive_proof_composition}
propõe composição recursiva de provas sem trusted setup. Vampire
\cite{vampire_zksnark} apresenta um zkSNARK universal. Inferência de
redes neurais com ZK é explorada em \cite{trustless_nn_inference}.

\lipsum[21]

% =============================================================================
% 3. PROBLEMA DE PESQUISA
% =============================================================================
\section{Problema de Pesquisa}

\lipsum[22-23]

\subsection{Hipótese}

\lipsum[24]

\subsection{Questões de Pesquisa}
\begin{enumerate}
  \item \lipsum[25][1-2]
  \item \lipsum[25][3-4]
  \item \lipsum[26][1-2]
\end{enumerate}

% =============================================================================
% 4. METODOLOGIA
% =============================================================================
\section{Metodologia Proposta}

\lipsum[27-28]

\subsection{Etapas da Pesquisa}
\begin{enumerate}
  \item \textbf{Revisão sistemática:} \lipsum[29][1-2]
  \item \textbf{Definição formal:} \lipsum[29][3-4]
  \item \textbf{Prototipação:} \lipsum[30][1-2]
  \item \textbf{Avaliação experimental:} \lipsum[30][3-4]
\end{enumerate}

\subsection{Ferramentas e Ambiente}

\lipsum[31]

% =============================================================================
% 5. CONTRIBUIÇÕES ESPERADAS
% =============================================================================
\section{Contribuições Esperadas}

\lipsum[32-33]

\begin{itemize}
  \item \lipsum[34][1-2]
  \item \lipsum[34][3-4]
  \item \lipsum[35][1-2]
\end{itemize}

% =============================================================================
% 6. CRONOGRAMA
% =============================================================================
\section{Cronograma (2 meses)}

\lipsum[36]

\begin{table}[h]
  \centering
  \begin{tabular}{|l|c|c|c|c|c|c|c|c|}
    \hline
    \textbf{Atividade} & \textbf{S1} & \textbf{S2} & \textbf{S3} & \textbf{S4} & \textbf{S5} & \textbf{S6} & \textbf{S7} & \textbf{S8} \\
    \hline
    Revisão da literatura        & $\bullet$ & $\bullet$ & $\bullet$ &           &           &           &           &           \\
    \hline
    Definição formal do problema &           & $\bullet$ & $\bullet$ & $\bullet$ &           &           &           &           \\
    \hline
    Desenvolvimento do modelo    &           &           & $\bullet$ & $\bullet$ & $\bullet$ & $\bullet$ &           &           \\
    \hline
    Prototipação e experimentos  &           &           &           & $\bullet$ & $\bullet$ & $\bullet$ & $\bullet$ &           \\
    \hline
    Escrita do pré-projeto       & $\bullet$ & $\bullet$ & $\bullet$ & $\bullet$ & $\bullet$ & $\bullet$ & $\bullet$ & $\bullet$ \\
    \hline
    Revisão e entrega            &           &           &           &           &           &           & $\bullet$ & $\bullet$ \\
    \hline
  \end{tabular}
  \caption{Cronograma de 8 semanas para o pré-projeto.}
\end{table}

% =============================================================================
% REFERÊNCIAS
% =============================================================================
\newpage
\bibliographystyle{abbrv}
\bibliography{references}

\end{document}
